\documentclass[11pt,a4paper]{article}

\usepackage[utf8]{inputenc}
\usepackage[T1]{fontenc}
\usepackage[portuguese]{babel}
\usepackage{amsmath,amsfonts,amssymb}
\usepackage{graphicx}
\usepackage{hyperref}
\usepackage[margin=1in]{geometry}
\usepackage{cite}

\title{\textbf{Métodos Escaláveis de Segunda Ordem e Híbridos para Otimização com Regularizações Descontínuas}\\
\large Projeto de Pesquisa - FAPESP BEPE}

\author{
    \textbf{Candidato:} Gabriel Belém Barbosa \\
    \textbf{Orientador Brasil:} Prof. Dr. Paulo José da Silva e Silva (UNICAMP) \\
    \textbf{Supervisor Exterior:} Prof. Dr. Dominique Orban (Polytechnique Montréal)
}
\date{Período Proposto: Setembro 2026 -- Agosto 2027}

\begin{document}

\maketitle

\begin{abstract}
Este é um projeto para estágio no exterior (BEPE) na Polytechnique Montréal sob a supervisão do Prof. Dominique Orban. O objetivo principal é desenvolver métodos escaláveis de segunda ordem e abordagens híbridas para problemas de otimização em larga escala que envolvem regularizações descontínuas, em especial a restrição de cardinalidade (norma $L_0$) e esparsidade de grupos. A proposta integra as heurísticas de primeira ordem (Nonmonotone Spectral Proximal Gradient - NSPG, descida de coordenadas e algoritmos combinatoriais locais) desenvolvidas durante o doutorado no Brasil com o maquinário avançado de otimização não-suave e métodos de confiança/regularização cúbica do grupo anfitrião. Além do desenvolvimento teórico e algorítmico, o projeto prevê a implementação robusta destas ferramentas no ecossistema de código aberto \texttt{JuliaSmoothOptimizers} (JSO), contribuindo diretamente para a comunidade de otimização.
\end{abstract}

\tableofcontents
\newpage

%LTeX: language=pt-BR
\section{Introdução e Justificativa}

A regularização $L_0$ e suas generalizações para esparsidade de grupos são ferramentas fundamentais em aprendizado de máquina e processamento de sinais, permitindo a seleção de variáveis e a interpretabilidade de modelos matemáticos complexos. Devido à natureza não-convexa e combinatória da penalidade $L_0$, a resolução eficiente desses problemas em larga escala permanece um desafio em aberto na otimização contínua. Outras regularizações, como híbridos de $L_1$ (Lasso) e $L_2$ (ridge) com $L_0$, também são de interesse recente por seu potencial para evitar overfitting, mantendo os benefícios da regularização pura $L_0$ \cite{fastselect}.

Durante o desenvolvimento do doutorado no Brasil, o foco central tem sido o desenvolvimento de heurísticas eficientes de primeira ordem para essas classes de problemas. Especificamente, desenvolvemos algoritmos baseados em Nonmonotone Spectral Proximal Gradient (NSPG) combinados com Descida de Coordenadas (CD) e métodos de busca local combinatória. Esses métodos não monótonos atuam de forma eficiente como um pré-processamento (identificação grosseira de suporte) para algoritmos do estado da arte (i.e., PGCCD \cite{fastselect}), superando desafios de dimensionalidade em ambientes de alta complexidade e correlação de paramêtros. A implementação foi desenvolvida utilizando a linguagem Julia, que permite alta performace e flexibilidade no desenvolvimento de algoritmos de otimização.

Apesar da alta eficiência computacional dos métodos de primeira ordem (espectrais), sua capacidade de atuar como pré-processadores globais e identificar boas regiões do espaço de busca poderia se beneficiar substancialmente de abordagens de segunda ordem. Aprimorar essa etapa com operadores proximais estruturados permite que o algoritmo se aproxime muito mais da variedade correta, garantindo que os métodos subsequentes de descida de coordenadas e algoritmos combinatórios tenham mais facilidade em realizar o refinamento local e encontrar soluções globalmente melhores. É neste contexto que a colaboração internacional proposta se fundamenta.

\subsection{Justificativa para a Escolha do Grupo Anfitrião}
O grupo de pesquisa liderado pelo Prof. Dominique Orban na Polytechnique Montréal (GERAD) possui liderança e excelência internacionalmente reconhecidas no desenvolvimento de métodos escaláveis de otimização não-suave e de segunda ordem. Trabalhos recentes do grupo, como o método de Levenberg-Marquardt para mínimos quadrados regularizados não-suaves \cite{aravkin2024levenberg} e métodos de regularização cúbica adaptativa escaláveis \cite{dussault2024scalable}, representam o estado da arte na resolução robusta de problemas complexos de otimização. Adicionalmente, análises recentes sobre a complexidade de métodos de região de confiança para otimização não-suave \cite{leconte2023complexity} oferecem uma base teórica sólida para extensões em regularizações descontínuas de interesse. Além disso, o grupo tem demonstrado forte interesse no estudo do comportamento algorítmico não monótono e no desenvolvimento de operadores proximais alternativos \cite{Leconte_2024}, o que se alinha perfeitamente com a trajetória de pesquisa atual do candidato.

Além do alinhamento teórico, o grupo hospedeiro é desenvolvedor principal do \texttt{JuliaSmoothOptimizers} (JSO), um ecossistema amplamente utilizado para otimização em Julia. A imersão neste ambiente de desenvolvimento propiciará uma sinergia única, permitindo não apenas o avanço teórico da pesquisa iniciada no Brasil, mas também sua conversão em software de alto impacto e uso global. Dessa forma, o estágio no exterior trará uma contribuição substancial que não poderia ser plenamente obtida na instituição de origem.

%LTeX: language=pt-BR
\section{Objetivos}

O objetivo principal deste estágio de pesquisa no exterior é estender e acelerar a resolução de problemas de otimização com regularizações descontínuas (tais como a penalidade $L_0$ e restrições de cardinalidade em grupos) por meio da exploração de novas alternativas de segunda ordem e operadores proximais estruturados. A intenção é utilizar esses métodos como poderosos pré-processadores globais capazes de navegar eficientemente por paisagens complexas e se aproximar da variedade ativa correta. Essa identificação aprimorada de boas regiões permitirá que refinadores locais --- como a descida de coordenadas e buscas combinatórias --- encontrem soluções de maior qualidade com muito mais facilidade.

Os objetivos específicos incluem:
\begin{itemize}
    \item \textbf{Desenvolvimento Algorítmico Híbrido:} Projetar um arcabouço algorítmico híbrido onde novas alternativas estruturadas de segunda ordem e operadores proximais (ex: métricas variáveis tridiagonais em blocos) atuem como identificadores globais de boas regiões. Isso será integrado à descida de coordenadas e às buscas locais combinatórias estudadas no doutorado, as quais atuarão como refinadores locais para isolar o suporte ótimo final;
    \item \textbf{Análise Teórica de Convergência:} Analisar a convergência local e global do método híbrido resultante, explorando as recentes análises de complexidade para otimização não-suave \cite{leconte2023complexity};
    \item \textbf{Estudo de Casos Práticos:} Validar os algoritmos propostos em aplicações reais de aprendizado de máquina esparso e processamento de sinais em grande escala;
    \item \textbf{Integração de Software:} Implementar os algoritmos resultantes em Julia, integrando as metodologias ao ecossistema \texttt{JuliaSmoothOptimizers} (JSO) se aplicável, garantindo alta escalabilidade e disponibilizando as ferramentas para a comunidade científica.
\end{itemize}

\subsection{Caminhos Alternativos de Pesquisa}
Embora o foco principal permaneça no arcabouço algorítmico híbrido, este projeto também identifica dois caminhos de pesquisa alternativos e promissores:
\begin{itemize}
    \item \textbf{Exploração de Separabilidade Inerente:} Aplicar o método de gradiente proximal com métrica variável \cite{Leconte_2024}, em conjunto com a descida de coordenadas e as buscas combinatórias, em problemas onde a Hessiana verdadeira é naturalmente parcialmente separável ou separável em blocos. Esta abordagem evita a necessidade de aproximações quasi-Newton que imponham tais critérios, aproveitando a estrutura nativa da função objetivo;
    \item \textbf{Técnicas Combinatórias para Otimização com Restrições:} Expandir as técnicas de busca combinatória para problemas com restrições (focando inicialmente em restrições de caixa). Neste contexto, a decisão de sair de uma face da região viável será tratada de forma análoga à decisão de adicionar ou remover parâmetros do suporte ativo em problemas $L_0$ sem restrições.
\end{itemize}


\section{Metodologia}

O projeto será desenvolvido em etapas incrementais e iterativas, garantindo o rigor metodológico característico da otimização contínua.

\subsection{Estudo Dirigido e Revisão Bibliográfica}
Inicialmente, será conduzido um estudo aprofundado dos métodos de região de confiança para otimização não-suave \cite{leconte2023complexity}, comportamento algorítmico não monótono e métodos de gradiente proximal indefinido \cite{Leconte_2024} desenvolvidos pelo grupo na Polytechnique Montréal. Paralelamente, haverá um treinamento focado na arquitetura do \texttt{JuliaSmoothOptimizers} para garantir que os desenvolvimentos posteriores estejam alinhados com as boas práticas de engenharia de software da comunidade.

\subsection{Desenvolvimento de Algoritmos Híbridos Ativos}
A metodologia de pesquisa focará na concepção de um modelo algorítmico híbrido. O componente desenvolvido no Brasil (NSPG + buscas combinatórias) atuará nas iterações iniciais do processo de otimização, promovendo esparsidade e identificando rapidamente um conjunto ativo (as variáveis não nulas). 

Adicionalmente, investigaremos o uso de matrizes de segunda ordem estruturadas no subproblema do operador proximal (operadores proximais de métrica variável). Embora abordagens recentes utilizando métricas diagonais indefinidas simples tenham apresentado resultados inferiores em comparação aos métodos espectrais padrão \cite{Leconte_2024}, a adoção de estruturas levemente mais densas, como matrizes tridiagonais, pode proporcionar um ganho significativo de convergência sobre o modelo espectral (múltiplo da identidade). De fato, nossos experimentos preliminares em mínimos quadrados com regularização L0 e estruturas tridiagonais já mostraram resultados promissores, demonstrando que tais subproblemas mantêm a rapidez computacional ao mesmo tempo em que oferecem alta precisão na recuperação do suporte.

Uma vez identificado o conjunto ativo pelos pré-processadores globais, o algoritmo aproveitará a descida de coordenadas e as buscas combinatórias para refinar estritamente a solução sobre as variáveis não nulas. Essa estratégia de \textit{active-set} reduzirá drasticamente a complexidade computacional, combinando a rápida navegação do espaço por operadores proximais estruturados com a convergência local precisa e isolamento de suporte dos refinadores CD.

\subsection{Avaliação Numérica e Benchmarking}
A implementação será feita na linguagem Julia, devido ao seu alto desempenho, que se alinha com o escopo do projeto no Brasil e a infraestrutura JSO no Canadá. Serão conduzidos experimentos extensivos comparando o método híbrido proposto contra algoritmos do estado da arte (como PGCCD e o pacote L0Learn) em instâncias sintéticas de alta dimensão e conjuntos de dados reais de aprendizado estatístico. As métricas de avaliação incluirão tempo de CPU, precisão de suporte, escalabilidade e qualidade da solução final em termos da função objetivo restrita.

\section{Plano de Trabalho e Cronograma}

O estágio de pesquisa terá duração de 12 meses (Setembro de 2026 a Agosto de 2027) e será dividido nas seguintes fases principais:

\begin{enumerate}
    \item \textbf{Mês 1-2: Integração e Estudo (Fase 1)} \\
    Instalação na instituição anfitriã; familiarização com a infraestrutura \texttt{JuliaSmoothOptimizers} (JSO); estudo de métodos de região de confiança e operadores proximais estruturados.
    \item \textbf{Mês 3-5: Concepção Algorítmica Híbrida (Fase 2)} \\
    Desenvolvimento teórico de operadores proximais estruturados como pré-processadores globais e sua integração com os refinadores locais de descida de coordenadas e combinatórios estudados no Brasil. Análise preliminar de viabilidade.
    \item \textbf{Mês 6-8: Implementação Computacional e Relatório Parcial (Fase 3)} \\
    Implementação paralela dos algoritmos projetados na linguagem Julia, integrando os modelos desenvolvidos no Brasil ao framework JSO. Elaboração do Relatório Científico anual da FAPESP (fevereiro de 2027).
    \item \textbf{Mês 9-10: Benchmarking e Validação (Fase 4)} \\
    Execução de baterias de testes escaláveis em problemas de regressão linear e logística com penalização $L_0$ e esparsidade em grupos. Comparação contra pacotes estabelecidos na literatura.
    \item \textbf{Mês 11-12: Consolidação e Redação (Fase 5)} \\
    Consolidação dos resultados, integração das metodologias desenvolvidas na tese de doutorado e elaboração do Relatório Científico Final do BEPE exigido pela FAPESP.
\end{enumerate}

O fluxo de transferência do conhecimento se dará por meio de reuniões regulares de acompanhamento online com o orientador no Brasil, garantindo que todo o desenvolvimento de código e avanços teóricos sejam diretamente incorporados à tese de doutorado.


\bibliographystyle{plain}
\bibliography{bibliografia}

\end{document}
